\section{Qwen3-Coder: Agentic Coding in the World}

La sección anterior examinó los Agent de producto de propósito general; esta sección examina los Agent en entornos de código. La pregunta central es por qué el coding es uno de los entornos de alto valor más adecuados para entrenar Agent y cómo los modelos de código abierto pueden alcanzar a los demás en este entorno. El tercer material es Qwen3-Coder. La publicación oficial del blog presenta Qwen3-Coder-480B-A35B-Instruct: MoE de 480B, 35B active, 256K de context nativo, ampliable a 1M mediante YaRN. Está orientado al agentic coding y destaca el code RL, el long-horizon RL, 20,000 entornos en paralelo y la cadena de herramientas Qwen Code.

\begin{figure}[H]
\centering
\includegraphics[width=0.96\textwidth]{qwen3-coder-stack.png}
\caption{Qwen3-Coder Stack: MoE de 480B, context de 256K/1M, Code RL y Long-horizon RL. Diagrama conceptual de elaboración propia, a partir de la publicación oficial del blog de Qwen y de la conversación de 01:53:38--01:59:04.}
\end{figure}

\begin{knowledgebox}{Lectura de la figura: el código es uno de los entornos más adecuados para el RL de Agent}
Las tareas de código son hard to solve, easy to verify: resolverlas es difícil, pero las pruebas y la retroalimentación de la ejecución son relativamente claras. Esto convierte al código en un buen entorno para el RL a gran escala y la long-horizon interaction.
\end{knowledgebox}

\subsection{Code RL y Long-horizon RL}

La subsección anterior presentó la pila del modelo Qwen3-Coder; esta explica por qué pone el énfasis en el RL. La publicación del blog de Qwen3-Coder subraya que las tareas de código se prestan de forma natural al execution-driven RL. Las tareas reales de ingeniería de software requieren interacción de varios turnos: planificar, usar herramientas, recibir retroalimentación y corregir decisiones. Qwen3-Coder introduce el long-horizon RL y construye entornos a gran escala que pueden ejecutarse en paralelo.

\begin{importantbox}{Por qué el coding es una posición estratégica para los Agent}
Los entornos de código ofrecen acciones claras, retroalimentación ejecutable, verificación mediante pruebas y tareas de alto valor. Frente a muchas tareas de mundo abierto, en el coding es más fácil definir la reward y también es más fácil escalar el entrenamiento y la evaluación.
\end{importantbox}

\begin{knowledgebox}{Por qué el Code RL es «difícil de resolver pero fácil de verificar»}
En muchas tareas de código el proceso de solución es difícil: exige entender la base de código, planificar los cambios y depurar; pero el resultado puede verificarse mediante pruebas, compilación, lint, benchmark o scripts del usuario. Esta estructura es muy adecuada para el aprendizaje por refuerzo, porque la reward puede automatizarse en buena medida.
\end{knowledgebox}

\subsection{Resumen de la sección}

Qwen3-Coder muestra la ruta sistemática de los modelos de código abierto en el agentic coding: contexto largo, datos de código, RL impulsado por la ejecución, interacción con herramientas a largo plazo y una cadena de herramientas de CLI.
